\section{Tracked files and history}
\label{sec:file-history}

A \emph{tracked file} carries its own provenance. Where \code{send}
sends a native file once, \code{keyquorum file} keeps a \file{.kqtf}
container that holds every revision of the file, the signatures over them,
the rules they are judged by, and an append-only, hash-chained history of
everything that happened to it. The container is the authority: SQLite only
keeps a cache of it, and a recipient can verify a tracked file without the
sender's database.

\subsection{What a tracked file holds}

\begin{itemize}
  \item \textbf{A stable identity.} A random 128-bit \code{file_id} is
  allocated once, when the file is first tracked, and never changes on
  rename, move, device transfer or new revisions. The logical file name is
  display metadata, not identity.
  \item \textbf{A revision graph.} Each revision records its parent
  revisions (none for the first, one for an edit, two for a merge), a
  content commitment bound to the file (so it is not a global fingerprint
  of the plaintext), the author's label and identity, a UTC time, the
  topology generation, and the hash of the file's policy. Its 32-byte
  \code{revision_id} is a domain-separated hash of all of that, so a
  revision cannot be re-parented or re-attributed without changing its id.
  Every revision also has a generated label,
  \code{R<file-name>-<UTC time>-<HCP label>}, and may have a user label;
  labels are for people, the id is what is signed.
  \item \textbf{Proofs.} Content signatures by a revision's author,
  countersignatures by an approver, and finalizations by the scope owner or
  an ancestor. They are ordinary Ed25519 signatures over domain-separated
  preimages (\code{KQ-FILE-REVISION-v1}, \code{KQ-FILE-COUNTERSIGN-v1} and
  \code{KQ-FILE-FINALIZE-v1}), made and checked by the same signing code as the
  rest of KeyQuorum.
  \item \textbf{A policy.} The HCP scope the file belongs to and what each
  kind of author needs before a revision is trusted (below).
  \item \textbf{A history.} Sequenced events, each linked to the hash of
  the one before; the last hash is the file's \emph{history root}. Changing,
  removing or reordering any earlier event breaks the chain. Events record
  outcomes and labels, never passwords, PINs, shares, bearer tokens or
  private keys. Every detail an event may carry comes from a fixed list of
  safe keys; anything else is refused before it is written. Some events are
  also signed by the person who caused them (below).
\end{itemize}

The container format is \code{KQTF}, currently version 7. Version 6 added
finalization proofs and version 7 optional event signatures; each is written
only when a file has one, so most files stay at an older version, and
versions 4 to 6 still open. It is not a sealed envelope and does not wrap or modify the
native file format: \code{file checkout} writes a revision's bytes back out
as an ordinary file.

\subsubsection*{When tracking starts, and what stays the same}

Signing a native file that has no container yet is the first content
signature: \code{file sign report.txt --as M.A --slot ...} allocates the
file id, records revision R1, signs it and writes
\file{report.txt.kqtf}, the same steps as \code{file track}. The file's
scope is the signer's own label unless \flag{--scope} names another. Files nobody
signs this way stay untracked. Tracking begins at the first \emph{trusted}
signature: if the first revision would still wait for a countersignature (a
descendant signing alone), nothing is written and the command says who can
start tracking; the explicit \code{file track} may still begin pending. A
second \code{file sign} on the same native path is refused instead of minting
a second identity. The check is by path: the same bytes signed from a copy at
another path start a separate tracked file, because a file's identity is a
random id and not a fingerprint of its content.

The file's identity does not depend on its name, path or holder.
\code{file rename} (scope owner or an ancestor, proven with \flag{--slot})
changes only the name shown, appends a \code{FILE_RENAMED} event with the
old and new names, and leaves the file id and every revision id alone.
Copying or moving the \file{.kqtf} file changes nothing either.
\code{file checkout} writes a revision's bytes back out exactly, whatever
the format.

\code{file status} also prints two ids that can differ, because a newer
edit is not a newer trusted version (and says when a revision is final):

\begin{itemize}
  \item \code{current_revision_id}: the sole head, or none when the history
  has forked.
  \item \code{trusted_revision_id}: the latest revision this store trusts.
\end{itemize}

Both are worked out on demand from the store's keys and are never stored in
the container, which carries no trust state.
They are views relative to the store that shows them, not claims the
container makes, so a copy of the file judged with other keys can show a
different trusted revision. On a fork there is no single current revision, and
the trusted one is only the most recently stored trusted revision across
both branches: a display value, never a branch winner and never the revision a
share would send.

\subsection{Trust: who may change a file}

A revision is \emph{trusted} only when its author's content signature
verifies against the signing key your store has registered for that label
\emph{and} the file's policy is met. The default policy (\code{file track})
is, relative to the file's scope:

\begin{center}
\begin{tabular}{@{}p{0.26\linewidth}p{0.68\linewidth}@{}}
\toprule
\textbf{Author} & \textbf{Needs} \\
\midrule
The scope owner & Their own signature. \\
A descendant (e.g.\ \code{M.A.1} in scope \code{M.A}) & Their own signature
  and a countersignature by their \emph{direct parent} (\code{M.A}). \\
An ancestor & Their own signature. \\
Cross-branch & Their own signature and either a private bridge's approval
  of that revision or a countersignature by the scope owner. \\
Unrelated & May not create trusted revisions. \\
\bottomrule
\end{tabular}
\end{center}

A revision is therefore \emph{pending} until its signatures are complete,
and \emph{denied} if a signature is present but wrong. A newer revision
never becomes trusted just because it is newer, and time stamps never decide
anything.

Each revision also records the topology generation (the version of the
published key tree) it was made under. Your store judges it only under a
generation it has held; a revision from a generation your store never saw is
left \emph{pending} (\code{MissingTopologyEvidence}) rather than judged
against today's tree. Generation 0 means no published tree.

The rules are chosen when the file is tracked and are fixed for its life.
Each role has its own option, which takes \code{forbidden}, \code{author}
or \code{author+parent}:

\begin{itemize}
  \item \flag{--owner-rule}
  \item \flag{--descendants-rule}
  \item \flag{--ancestors-rule}
  \item \flag{--cross-branch-rule}, which may also take
  \code{author+bridge-or-owner} or the stricter
  \code{author+bridge+owner}.
\end{itemize}

The defaults are the table above. A rule that cannot be met --- a forbidden
owner, a bridge for anyone but a cross-branch author, a parent
countersignature where there is no parent --- is refused before anything is
written. \code{file policy} prints the rules and the \code{policy_hash} that
every revision carries, so later verification knows which rules applied. When
a countersignature arrives, the history records its \code{relationship} to
the author (for example \code{DIRECT_PARENT}) and the policy decision records
its \code{reason}.

Revision labels are \code{R<name>-<UTC time>-<HCP label>} with the time to the
millisecond where the clock has it, so adjacent check-ins get different
labels. A label is part of the immutable revision: changing it would mean a
different revision, never an edit of the existing one.

\begin{lstlisting}
keyquorum file track report.txt --scope M.A --as M.A --slot ./usb=M.A
keyquorum file checkin report.txt.kqtf --from edited.txt --as M.A.1 --unsigned
keyquorum file sign report.txt.kqtf --as M.A.1 --slot ./usb=M.A.1
keyquorum file countersign report.txt.kqtf --as M.A --slot ./usb=M.A
keyquorum file status report.txt.kqtf
\end{lstlisting}

\subsubsection*{Bridge approval}

A cross-branch author can be approved through a private bridge instead of by
the scope owner. A member of a live private bridge whose members reach from
the author to the file's scope signs an approval of that exact revision:

\begin{lstlisting}
keyquorum file bridge-approve report.txt.kqtf --bridge <uid> --as M.S.1 \
  --slot ./usb=M.S.1
\end{lstlisting}

The approval is kept only in the approving store and never written into the
\file{.kqtf}, so the file never names the bridge or its members. Each time a
revision is judged it is checked again: the bridge must still be live at the
same generation and the signer still a member. A bridge's supervisors do not
count, and a bridge merely existing between two labels approves nothing. A
recipient whose store holds no such approval sees the revision as pending.
Under \code{author+bridge-or-owner}, a valid scope-owner countersignature
then substitutes for the bridge approval. A trusted cross-branch policy
decision records the alternative that satisfied it as
\code{satisfied_by=BRIDGE}, \code{SCOPE_OWNER}, or
\code{BRIDGE_AND_SCOPE_OWNER}. It never records a bridge id, generation,
signer, member, or roster, but \code{BRIDGE} and
\code{BRIDGE_AND_SCOPE_OWNER} do reveal that a bridge approval took part in
the decision, and that stays in the \file{.kqtf} and in any export of it.

With the stricter \code{author+bridge+owner} rule, the scope owner's
countersignature does not replace the bridge approval. The receiving store
must hold and validate that revision-specific bridge approval \emph{and} the
container must carry the scope owner's valid countersignature before the
revision is trusted.

\subsubsection*{Finalization}

Trusted is not final. The scope owner or one of its ancestors can mark a
trusted revision final:

\begin{lstlisting}
keyquorum file finalize report.txt.kqtf --as M.A --slot ./usb=M.A
\end{lstlisting}

This adds a signed finalization proof and a \code{REVISION_FINALIZED} event.
Like trust, finality is worked out by each store from its own keys; a store
without the finalizer's key does not see the revision as final. A merge of
two final revisions is a new revision and starts neither trusted nor final.

\begin{securityBox}
Your store supplies only public signing keys, and it accepts a key only for
the identity it knows for that label (the one enrolled by
\code{transfer enroll}, or else one derived from the label). A container
cannot bring its own keys. A ghost identity (after a \code{transfer move})
cannot author or sign anything new; what it signed earlier still verifies.
\end{securityBox}

\subsection{Sharing falls back to the last trusted revision}

\code{send} (which runs \code{file share} for a tracked file) sends the head
(or the revision named with \flag{--revision}) when it is \emph{trusted}. If it is not --- say someone
checked in an edit without signing it --- the letter carries its nearest
trusted \emph{ancestor} instead. Trusted work on another branch is never
substituted, and the unsigned revision never leaves your copy: the container
in the letter holds only the delivered revision and its ancestors. If no
ancestor is trusted, nothing is sent and the refusal is recorded. The
refusal says why: \emph{denied by policy} when the file's rules refuse the
revision outright (an unrelated or forbidden author, evidence that does not
verify), and \emph{no trusted revision} when it is only waiting for a
signature.

\begin{lstlisting}
# M.A (identity and relay stored with `use` and `loadkey`)
keyquorum send report.txt.kqtf --to M.B
# M.B: list what is waiting, then open the letter into a new copy
keyquorum inbox
keyquorum inbox open <id> --out report.txt.kqtf
# M.A, later: record M.B's signed answer in M.A's copy
keyquorum inbox open <answer-id> --file report.txt.kqtf
\end{lstlisting}

Letters are \code{KQPB} envelopes of their own kind, carried by the relay
like any other letter. The sender's signature on the letter is
\emph{transport} authentication only: the recipient still judges the
delivered revision by the file's own policy against their own store, and
refuses it if it is not trusted there. The recipient's store must also hold the encryption key that
opened the letter under the recipient label the letter names, so a letter
sealed to one person cannot credit the delivery to another; \code{inbox open} (and the older
\code{deliver open} and \code{file receive}) make the same check. The signed
letter also names every proof it carries for the delivered revision, and the
receiving command refuses a letter whose container does not match. A rename made before sharing travels with the
letter, and so does any later history that does not concern a revision
outside the delivered one. \flag{--into} merges a letter into a
copy you already hold. Either way the recipient seals back a signed accept
or reject, and recording it (\code{inbox open <id> --file}, or the older
\code{file ack}) happens once in the sender's history.
Receiving the same letter again records nothing new; it only seals the
answer again.

Receiving also says how the delivered history compares with what
your store accepted before for that file: \code{FIRST}, \code{REPLAYED} (the
same revision and history again), \code{NEWER} (its history passes through
every history root you accepted before, holds those revisions and adds one)
or \code{NOT_NEWER} (an older copy, another branch, or the same revisions on
a rewritten history). It is recorded on the delivery. \code{NOT_NEWER} is a
warning, not a refusal: a trusted older revision is still accepted.

\code{file send-history} seals a file's event history alone (a \code{KQHS}
snapshot, no content) to another label, with its root and event count signed
by you. Opening it (\code{inbox open <id> --file COPY}, or the older
\code{file open-history --against}) checks the signature and the snapshot and
compares it with your copy: \code{SAME}, \code{LOCAL_AHEAD},
\code{REMOTE_AHEAD} or \code{DIVERGED}. It is not answered.

\begin{lstlisting}
keyquorum file send-history report.txt.kqtf --to M.B --as M.A --push
# M.B:
keyquorum inbox open <id> --file report.txt.kqtf
\end{lstlisting}

\subsection{Sending and receiving a tracked file from the command line}

A tracked file travels as a \emph{letter}: a sealed \file{.kqpb} file that
holds no secret. It is signed by the sender and sealed to the recipient's
public key, so it can cross any channel. With a relay the commands are
\code{send} and \code{inbox open}; without one the letter and the answer are
files you carry, and only the way they get from one person to the other
differs. Each person needs only their \emph{own} drive, plugged into their \emph{own} computer,
and only at the moment they sign or open something. Nobody hands over a
key drive.

The examples send \file{report.txt.kqtf} from M.A to M.B. Each command names
the person who runs it and the computer it runs on.

\subsubsection*{Step 0, once, before either way: exchange public keys}

Each store needs the other's public keys: the sender's signing key (to check
the letter) and the recipient's encryption key (to seal it). They are
public, so any channel will do.

\begin{lstlisting}
# 0a. on M.A's computer (M.A's drive in): print M.A's public keys, send them to M.B
keyquorum-device public ./usb --label M.A
# 0b. on M.B's computer: register what M.A sent
keyquorum --db mb.sqlite register --type encryption --label M.A \
  --public-key-file a-enc.pub
keyquorum --db mb.sqlite register --type signing --label M.A \
  --public-key-file a-sign.pub
# 0c. M.B does the same in the other direction: prints their keys with
#     keyquorum-device public, M.A registers them in ma.sqlite
\end{lstlisting}

\subsubsection*{With a relay: send, then inbox open}

A relay is a mailbox that stores sealed letters and cannot open them. It
adds no security the letter lacks. It saves the two people from handing
the files over and lets them act at different times. You need the relay's
URL and API keys from its operator. Each key has one scope: an
\emph{inbox push} key sends a letter or an answer, an \emph{inbox pull} key
downloads what is waiting, so each person usually holds both. Store the
identity and the relay once per person, on their own computer (\code{setup},
or \code{use} and \code{loadkey}); after that the commands need no
\flag{--as}, \flag{--slot} or \flag{--url}.

\begin{lstlisting}
# 1. M.A's computer, M.A's drive in: seal the letter and upload it
keyquorum send report.txt.kqtf --to M.B

# 2. Later, M.B's computer, M.B's drive in: see what is waiting (id, kind)
keyquorum inbox

# 3. M.B opens the letter and uploads the signed answer
keyquorum inbox open <id> --out report.txt.kqtf
#    --into report.txt.kqtf   merge into a copy M.B already holds (instead of --out)
#    --reject                 refuse the letter (instead of --out or --into)

# 4. Later, M.A's computer, M.A's drive in: record the answer in M.A's copy
keyquorum inbox open <answer-id> --file report.txt.kqtf
\end{lstlisting}

The copy a letter concerns is always named by the person
(\flag{--into}, \flag{--out} or \flag{--file}), never taken from the letter:
its file id is only the sender's claim, and a copy that does not match is
refused. \flag{--url} and \flag{--api-key} can be given on \code{send} and
\code{inbox} instead of using the stored key. These environment variables work
too:

\begin{itemize}[leftmargin=1.5em]
  \item \code{KEYQUORUM_RELAY_URL}
  \item \code{KEYQUORUM_RELAY_API_KEY}
\end{itemize}

Downloaded letters are kept under \file{./inbox/<relay hash>/<n>.kqpb}, named
by the relay's own number, so letters from two relays can never be mistaken
for each other. The official client only talks to a relay that proves it holds
a KeyQuorum-signed certificate. Everything else, including the sender's
signature, the recipient-key check, the history comparison and the signed
answer, is the same as without the relay.

\subsubsection*{Without a relay: you carry the files}

No URL and no API key. The letter and the answer are plain files; move them
by email, a shared folder, or a USB stick that is only a courier.
\code{inbox} reads a relay, so a letter that arrives as a file is opened with
\code{file receive}, and its answer is recorded with \code{file ack}. Both
stay supported for exactly this case.

\begin{lstlisting}
# 1. M.A's computer, M.A's drive in: seal the letter into a folder
keyquorum --db ma.sqlite send report.txt.kqtf --to M.B --as M.A \
  --slot ./usb=M.A --output-dir ./outbox

# 2. Carry ./outbox/<id>.kqpb to M.B's computer (any channel)

# 3. M.B's computer, M.B's drive in: open the letter and seal the answer
keyquorum --db mb.sqlite file receive --letter ./<id>.kqpb \
  --slot ./usb=M.B --ack-dir ./acks --out report.txt.kqtf
#    --into report.txt.kqtf   merge into a copy M.B already holds (instead of --out)
#    --reject                 refuse the letter (instead of --out or --into)

# 4. Carry ./acks/<id>-ack.kqpb back to M.A's computer (any channel)

# 5. M.A's computer, M.A's drive in: record the answer
keyquorum --db ma.sqlite file ack report.txt.kqtf --ack ./<id>-ack.kqpb \
  --slot ./usb=M.A
\end{lstlisting}

To send M.B's changes back, swap the roles and repeat. If both copies have
moved on, the receiver ends up with two heads:

\begin{lstlisting}
# 6. on whoever received: join the heads, then settle any conflict
keyquorum --db ma.sqlite file merge report.txt.kqtf --as M.A
keyquorum --db ma.sqlite file review report.txt.kqtf
keyquorum --db ma.sqlite file resolve report.txt.kqtf --keep right \
  --as M.A --slot ./usb=M.A
keyquorum --db ma.sqlite file sign report.txt.kqtf --as M.A --slot ./usb=M.A
\end{lstlisting}

A no-drive alternative to steps 1 to 5 is to carry the whole \file{.kqtf}
and run \code{file import report.txt.kqtf --from theirs.kqtf --as M.B}. It
skips the signature and trust check that a letter carries.

\subsubsection*{Legacy procedure with a relay (still possible, not recommended)}

Before \code{send} and \code{inbox open}, the relay path was five separate
commands, with a download and an upload step the person had to run by hand.
They still work as they always did, and each prints a note on standard error
naming the command that replaces it. Prefer the commands above.

\begin{lstlisting}
keyquorum --db ma.sqlite loadkey --url https://relay.example
keyquorum --db ma.sqlite file share report.txt.kqtf --to M.B --as M.A \
  --slot ./usb=M.A --push                                  # M.A uploads
keyquorum --db mb.sqlite relay pull --output-dir ./inbox   # M.B downloads
keyquorum --db mb.sqlite file receive --letter ./inbox/<n>.kqpb \
  --slot ./usb=M.B --push-ack --out report.txt.kqtf        # M.B answers
keyquorum --db ma.sqlite relay pull --output-dir ./answers # M.A downloads
keyquorum --db ma.sqlite file ack report.txt.kqtf --ack ./answers/<n>.kqpb \
  --slot ./usb=M.A                                         # M.A records
\end{lstlisting}

In that form, letters downloaded by \code{relay pull} are named by the relay's
number only (\code{<n>.kqpb}) and \flag{--url} and \flag{--api-key} go on
\code{share} and \code{receive}.

\begin{cautionBox}
Without the relay a drive is needed only for the person who is signing or
opening at that moment. The only cases that need two devices open together
are moving or copying an identity between devices
(\code{keyquorum transfer copy|move}), and the browser lab, which simulates
everyone on one machine. A merge proposal, a file request and a change
request all travel as letters the same way (see the two sections below).
\end{cautionBox}

\subsection{Copies, forks and merging}

Two holders of a file can each extend it. \code{file import} brings another
copy's revisions and proofs into yours: nothing is overwritten, and if both
sides added revisions the file now has two heads. \code{file merge} then
tries a conservative automatic merge of UTF-8 text, line by line, against
the heads' nearest common ancestor:

\begin{center}
\begin{tabular}{@{}p{0.26\linewidth}p{0.68\linewidth}@{}}
\toprule
\textbf{Outcome} & \textbf{Meaning} \\
\midrule
\code{FastForward} & One head already contains the other. \\
\code{AlreadyEquivalent} & Both heads hold the same content. \\
\code{CleanMerge} & The edits do not touch the same lines; a new
  two-parent revision is recorded. \\
\code{RequiresHuman} & The edits overlap differently, or there is no single
  common ancestor. \\
\code{PolicyBlocked} & The file's policy turns automatic merging off
  (\flag{track --no-auto-merge}). \\
\code{UnsupportedContent} & The content is not UTF-8 text. \\
\bottomrule
\end{tabular}
\end{center}

A merge revision never inherits its parents' signatures: it starts
\emph{pending} and has to be signed (and countersigned where the policy
says so) like any other revision. When a person has to decide,
\code{file merge} records the conflict and names a reviewer: the most
recent neutral prior owner of the file, else the most senior eligible
owner, else (for a tie across sectors) the larger sector, else the nearest
common senior ancestor. Only owners whose own revisions are trusted
qualify, the authors of the conflicting edits never review their own
collision, and if nobody qualifies the conflict is recorded as unresolved
rather than guessed; then the scope owner or one of its ancestors decides.

The reviewer, or one of their ancestors, settles the conflict with
\code{file resolve}: keep one side as it is (\flag{--keep left} or
\flag{--keep right}) or supply an edited result (\flag{--from}). The result
is a new two-parent revision signed as the reviewer's own content and
recorded as \code{CONFLICT_RESOLVED}; the file's rules then judge it like any
other revision. \flag{--reject} instead refuses an untrusted proposed merge
at the head (\code{MERGE_REJECTED}): the proposal stays in the history,
\code{file sign} refuses it, and the resolution builds on it. A conflicting
author can never decide.

\begin{lstlisting}
keyquorum file resolve report.txt.kqtf --keep left --as M.A --slot ./usb=M.A
keyquorum file resolve report.txt.kqtf --from merged.txt --as M.A \
  --slot ./usb=M.A
keyquorum file resolve report.txt.kqtf --reject --as M.A --slot ./usb=M.A
\end{lstlisting}

\begin{lstlisting}
keyquorum file import report.txt.kqtf --from their-copy.kqtf --as M.A
keyquorum file merge report.txt.kqtf --as M.A
keyquorum file review report.txt.kqtf
keyquorum file diff report.txt.kqtf
keyquorum file graph report.txt.kqtf
\end{lstlisting}

\code{file review} prints each side's changed lines with who wrote each
line, the merge outcome and the reviewer. When the only head is a merge
that waits to be signed, \code{file review} shows what each parent changed
and what the merge changed (\code{LEFT PARENT}, \code{RIGHT PARENT},
\code{MERGED}) with its trust state, and \code{file merge} opens the
interactive review for it. Revisions are shown by their user
label first, when they have one, then the generated label and the short id.
A build with \flag{--features tui} adds \flag{review --interactive}: the same
view in the terminal, with removed text on a red and added text on a green background, Vim-style keys and mouse hover to see who wrote a
line. Press \code{za}, \code{zc} or \code{zo} to fold, close or open the change under the cursor
(a folded change shows its first line and a count), and \code{zM} or \code{zR} to fold or open every change on the side.
Two unchanged lines are shown on each side of a change, dimmed and marked as unchanged, and \code{c} hides or shows them; a fold hides them with the change, and the printed \code{file review} lists changed lines only. Press \code{Space} to pick the change under the cursor on either side,
\code{:compose} to see the result (the common ancestor with the picked changes applied;
picks that touch the same lines are refused), and \code{:edit} to change that
result in \code{$VISUAL} or \code{$EDITOR}, in an owner-only temporary file that is
removed afterwards. \code{:as LABEL} and \code{:slot CONTAINER=LABEL} (or
\flag{--as} and \flag{--slot} on \code{review --interactive}) say who is acting;
after that \code{:accept left}, \code{:accept right}, \code{:accept result},
\code{:reject}, \code{:sign} and \code{:finalize} run the real \code{file resolve},
\code{file sign} and \code{file finalize} with the terminal handed back for the
passphrase, so the reviewer and scope rules are those commands', and the review
then shows the history as it stands. \code{:accept result} settles the conflict when one waits for its reviewer; on a clean merge, which has no conflict, it checks the result in as a new signed revision (\code{file checkin}) on top of the merge. Press \code{/} to type a query and \code{Enter} to search; \code{n} and
\code{N} repeat it forward and backward. Matching ignores case, looks for a
substring in the active pane, and wraps around. On such a build,
\code{file merge} opens it by itself when a person
has to decide and the terminal is interactive. The review decides no trust
and edits no history itself: every action above is the real command run
as the label and slot you gave, and it is refused when that command would
refuse you. Without \code{:as} and \code{:slot} the commands only name the
CLI command to run.

\begin{cautionBox}
Commands replace a container under an exclusive \file{<name>.kqtf.lock}, and
refuse to write if the container changed since the command read it; run the
command again. If a command was killed and left the lock behind, the error
names the lock file to remove.
\end{cautionBox}

\subsection{Terms for proposing a change}

Three words describe someone asking for, or offering, a change to a file.
Each names something this version does. A request only asks: it delivers
nothing and changes nothing, and the answer to it only says yes or no. The
file or the change itself always travels the ordinary way, as a letter
(\code{send}) whose revision is judged by the file's own rules.

\begin{center}
\begin{tabular}{@{}p{0.22\linewidth}p{0.72\linewidth}@{}}
\toprule
\textbf{Term} & \textbf{Meaning and how it works} \\
\midrule
\textbf{Merge proposal} & Prepared changes submitted for review: a merge
  revision, made automatically from two heads or by an author outside the
  file's scope, that stays \code{Pending} until the reviewer, the scope
  owner or a bridge decides. The reviewer settles it with \code{file
  resolve} or refuses it with \flag{--reject}. The manual also calls it a
  proposed merge. \\
\textbf{Change request} & Asking someone to change a file (``please change
  this file''), before any change exists: \code{file request --change} with
  a short message. The holder reads it and answers with \code{inbox open ID
  --accept} or \flag{--decline} (offline: \code{file open-request} and
  \code{file answer-request}). \\
\textbf{File request} & Asking the holder to send a file, usually one you do
  not hold: \code{file request}. An accepted file request is served with
  \code{send}. \\
\bottomrule
\end{tabular}
\end{center}

Use \emph{change request} for ``please change this file'' and \emph{merge
proposal} once the author has already prepared the changes.

\subsection{Asking for a file, or for a change}

A request is a small signed letter (\file{.kqpb}) sealed to the holder, and
the answer is another, sealed back to the requester. They travel exactly as
in the previous section: through the relay, or as plain files you carry. The
holder checks the requester's signature against the key their own store
holds for that label, and the requester checks the holder's answer the same
way. A request never opens as a delivery letter, or the other way round.

What a history records, when a copy of the file is named: the request
(\code{FILE_REQUESTED} or \code{CHANGE_REQUESTED}) with its id and kind, and
the answer (\code{REQUEST_ANSWERED}) with the decision. The message of a
change request is shown to the holder and is \emph{never} recorded. Opening
the same request or answer twice records it once, and a request already
answered one way cannot be answered the other.

\subsubsection*{With a relay: file request, then inbox open}

Store the relay URL and key once per person (\code{loadkey}, or \code{setup}).
M.B asks M.A; each answer is one \code{inbox open}.

\begin{lstlisting}
# 1. M.B's computer, M.B's drive in: ask for a file M.B does not hold yet
keyquorum file request --file-id <32 hex digits> \
  --name report.txt --to M.A --as M.B --push
#    or ask for a change to M.B's own copy (message up to 1024 bytes)
keyquorum file request report.txt.kqtf --change \
  --message "the total in row 3 is wrong" --to M.A --as M.B --push

# 2. Later, M.A's computer, M.A's drive in: see what is waiting
keyquorum inbox

# 3. M.A reads the request and answers it in one command; --file records it
#    in M.A's copy, and a decision is never guessed
keyquorum inbox open <id> --accept --file report.txt.kqtf   # or --decline

# 4. Later, M.B's computer. A file request: read the answer by its id (M.B
#    holds no copy to record it in)
keyquorum inbox open <answer-id>
#    A change request: record the answer in the copy that recorded the request
keyquorum inbox open <answer-id> --file report.txt.kqtf
\end{lstlisting}

A request that is waiting is listed by \code{inbox} and left alone by a bare
\code{inbox open}, because it needs a decision only M.A can give.

If M.A accepted a \emph{file request}, M.A now sends the file with
\code{send}, as in the previous section. If M.A accepted a \emph{change
request}, the change reaches M.A as a revision the ordinary way: M.B checks it
in on their copy and sends it with \code{send}, and it arrives as a merge
proposal.

\subsubsection*{Without a relay: you carry the files}

The public keys of both people are exchanged and registered first (step 0
in the previous section). \code{inbox} reads a relay, so a request or an answer
that arrives as a file is opened with the \code{file} commands, which stay
supported for this case.

\begin{lstlisting}
# 1. M.B's computer, M.B's drive in: ask for a file M.B does not hold yet
keyquorum --db mb.sqlite file request --file-id <32 hex digits> \
  --name report.txt --to M.A --as M.B --slot ./usb=M.B --output-dir ./requests
#    or ask for a change to M.B's own copy (message up to 1024 bytes)
keyquorum --db mb.sqlite file request report.txt.kqtf --change \
  --message "the total in row 3 is wrong" --to M.A --as M.B \
  --slot ./usb=M.B --output-dir ./requests

# 2. Carry ./requests/<id>-request.kqpb to M.A (any channel)

# 3. M.A's computer, M.A's drive in: read it; --file records it in M.A's copy
keyquorum --db ma.sqlite file open-request --letter ./<id>-request.kqpb \
  --slot ./usb=M.A --file report.txt.kqtf

# 4. M.A's computer: answer (accept or decline); --file records the answer
keyquorum --db ma.sqlite file answer-request --letter ./<id>-request.kqpb \
  --decision accept --slot ./usb=M.A --file report.txt.kqtf --ack-dir ./answers

# 5. Carry ./answers/<id>-answer.kqpb back to M.B (any channel)

# 6. M.B's computer, M.B's drive in: read the answer (--file records it
#    against the request M.B's copy shows it made)
keyquorum --db mb.sqlite file open-answer --answer ./<id>-answer.kqpb \
  --slot ./usb=M.B --file report.txt.kqtf
\end{lstlisting}

\subsubsection*{Older procedure with a relay (still possible, not recommended)}

The same exchange used to be a download and a separate command at every step.
These \code{file} commands print no note, but they still work, and \code{relay pull --import} still lists a request or an
answer and tells you which command opens it; it never imports one, because a
request only asks. Prefer \code{inbox open} above.

\begin{lstlisting}
keyquorum --db ma.sqlite relay pull --output-dir ./inbox
keyquorum --db ma.sqlite file open-request --letter ./inbox/<n>.kqpb \
  --slot ./usb=M.A --file report.txt.kqtf
keyquorum --db ma.sqlite file answer-request --letter ./inbox/<n>.kqpb \
  --decision accept --slot ./usb=M.A --file report.txt.kqtf --push-ack
keyquorum --db mb.sqlite relay pull --output-dir ./answers
keyquorum --db mb.sqlite file open-answer --answer ./answers/<n>.kqpb \
  --slot ./usb=M.B --file report.txt.kqtf
\end{lstlisting}

\subsection{History snapshots and the index}

\code{file history export FILE --out SNAPSHOT} writes a \code{KQHS} snapshot: the event
history alone, which verifies on its own. \code{file history verify SNAPSHOT}
with \flag{--against FILE} checks that a snapshot is a point in a file's history --- for
example, what a partner held when they last answered. A snapshot of an older
point is still valid; it proves a prefix, not that nothing came after. The older spellings
\code{file history --export} and \code{file verify-snapshot} still work but are
not recommended.

The store keeps \code{tracked_files}, \code{tracked_revisions} and
\code{tracked_history_index} as a cache of metadata (no content, no trust
decisions). \code{file list} reads it and \code{file reindex} rebuilds it
from the containers, which remain the authority.

\subsection{Expiry and tombstones}

\code{file expire} lets the scope owner (or one of its ancestors) schedule
when a tracked file's content ends, or end it now. Because this cannot be
undone, naming the owner is not enough: the command signs a challenge with
\flag{--slot} (or \flag{--signing-key-file}) and checks it against the key
your store has registered for that label. The key is opened only for that
proof and released before anything changes. When the time has passed,
the next command that needs the content destroys \emph{every} retained
revision's payload at once and records \code{FILE_EXPIRED} and
\code{CONTENT_DESTROYED}. What remains is a tombstone: the revision graph,
signatures and history still verify, no new revision can be added, and each
later attempt to use the content is recorded as
\code{EXPIRED_ACCESS_ATTEMPT} and refused. A container that is missing some
payloads without a recorded destruction, or keeps any after one, fails
verification.

\begin{lstlisting}
keyquorum file expire report.txt.kqtf --as M.A --at "2026-12-31 17:00" \
  --slot ./usb=M.A
keyquorum file expire report.txt.kqtf --as M.A --now --slot ./usb=M.A
\end{lstlisting}

\begin{cautionBox}
Expiry destroys the content of \emph{this} container. Copies someone
already holds --- a letter they accepted, a checkout they wrote out --- are
their own files and are not reached.
\end{cautionBox}

\subsection{Signed events}

The hash chain shows that no earlier event was changed, removed or
reordered. Commands that already prove you hold your key --- \code{file
rename}, \code{file expire} and \code{file resolve --reject} --- also sign
the events they append (over \code{KQ-FILE-HISTORY-EVENT-v1}). \code{file history} marks such an event
\code{[signed by ...]} where your store's key for that person verifies it,
and \code{[signature not verified in this store]} where it does not. Other
events are protected by the chain alone and are never shown as signed. An
event signature says who recorded the event; it decides nothing about
whether a revision is trusted.

\subsection{Recording what happens at a gate}

\code{file link} ties a quorum-protected file (\flag{--quorum-file}) or a
password-locked file (\flag{--locked-file}) to a tracked file. From then on
the CLI appends what happened at that gate to the tracked file's history:
each unlock attempt and its outcome (with the labels presented on a quorum
success), each share link created, redeemed or revoked (the share's id,
never its token), and the tombstone when the gate's own expiry destroys its
ciphertext. The gates themselves are unchanged and never read the link;
recording happens after the gate has answered and cannot change it. An
unlock that fails before reaching the gate (for example while shares are
being collected) is not recorded.

What is recorded is limited to what is safe to keep:

\begin{itemize}
  \item \textbf{Quorum unlocks:} counts and policy words only --- shares
  presented, threshold, distinct devices, minimum devices, custody mode,
  approval mode and parent approvals --- never a share or key.
  \item \textbf{PIN checks:} only the outcome (\code{asked},
  \code{verified}, \code{mismatch}, \code{locked} or \code{failed}); nothing
  when no PIN applies. Never the PIN, its hash or the attempt count.
  \item \textbf{Share links:} the share's id, never its token. A redemption
  is recorded as \code{redeemer=UNKNOWN_BEARER}, because holding a token
  proves nothing about who is holding it.
\end{itemize}

\subsection{Limits this version states}

Two limits are stated on purpose, so this version does not claim stronger
guarantees than it has.

\begin{itemize}[leftmargin=1.5em]
  \item \textbf{Historical authority is local key history.} The store
  preserves which identity and signing key each label held, as observed
  here. It does not establish complete generation-specific identity,
  relationship or reviewer authority.
  \item \textbf{Bridge disclosure awaits the owner's acceptance.} A
  \code{satisfied_by} of \code{BRIDGE} or \code{BRIDGE_AND_SCOPE_OWNER}
  reveals that a bridge approval took part in a decision. That stays in the
  portable history and in any export. Whether to keep it is the owner's
  policy decision.
\end{itemize}

The full architecture goes further than what is built. Today:

\begin{itemize}
  \item Tracking starts with \code{file track} or the first \code{file sign}
  of a native file; nothing tracks a file on its own. A request for a file
  or a change is recorded only in a copy you name, and never as a delivery.
  \item A store keeps which topology generations it has held. Each command that
  resolves a file's topology generation (track, check in, sign, merge and
  the like, not a plain \code{verify}) also records which identity and
  signing key each label under that tree holds now, as the first and last
  generation it was observed in. A revision stamped with an older
  generation is also checked against every key recorded for that label and
  identity whose observed range covers it, in addition to the current key,
  so it keeps its meaning after the label is reissued or reassigned. This is
  the store's own key history, not a per-generation record of who held the
  label: a generation between two observations counts as covered, the
  current key is tried without proof it held the label then, a store that
  never saw a key has no record of it, and nothing proves when a signature
  was made.
  \item Conflicts are settled a whole revision at a time (\flag{--keep} or
  \flag{--from}), not hunk by hunk. The interactive review settles a conflict
  from the terminal only through the same commands as the command line, and
  a line has one recorded author in it: where several revisions on one side
  touched it, only the last is shown.
  \item Checkouts, signature checks and exports are recorded only when asked
  (\flag{--record}), and only under the label \flag{--as} names; that label is
  a claim, so the event is chained but never attested. Tampering is refused
  when a file is opened rather than recorded in it.
  \item Recording at gates is opt-in (\code{file link}) and best effort.
\end{itemize}
